\chapter{LoRA and Full Fine-Tuning: Accuracy and Training Cost}
\label{ch:lora_ft}

This chapter compares LoRA on all blocks (rank 8) with full fine-tuning on both models.
Sections~\ref{sec:lr} to~\ref{sec:targets} are about accuracy, Sections~\ref{sec:mem} to~\ref{sec:size} about the
cost of training.

\section{Choosing the learning rates}
\label{sec:lr}

A comparison of two training methods is only fair if each is run at a learning rate that suits it. Full
fine-tuning changes all pretrained weights and needs a small rate; LoRA trains small matrices that start at zero
and tolerates a larger one. We therefore ran each method at four rates on each model, with three seeds (42, 123,
456), and chose the rate with the lowest validation loss (Table~\ref{tab:lr}, Figure~\ref{fig:lr}).

\begin{table}[!htb]
\centering\small
\begin{tabular}{llccc}
\toprule
method & learning rate & validation loss & test acc.\ (\%) & stopping epoch \\
\midrule
\rawinput{data/tab_lr.tex}
\bottomrule
\end{tabular}
\caption{Full fine-tuning and attention LoRA (rank 8, all blocks) at four learning rates each; mean $\pm$ standard
deviation over 3 seeds. $^{*}$ = lowest mean validation loss of the method on that model. The ViT-B runs of full
fine-tuning at $10^{-4}$ are from an earlier set of runs with the same settings.}
\label{tab:lr}
\end{table}

\begin{figure}[!htb]
\centering
% generated by scripts/make_figures.py; do not edit
\begin{tikzpicture}
\begin{groupplot}[group style={group size=2 by 2, horizontal sep=\figSep, vertical sep=0.5cm}, width=\figW, height=\figH, repaxis, xmode=log, xmin=1.8e-6, xmax=1.7e-3, ymajorgrids, legend style={font=\figLegendFont, draw=none, fill=none, at={(-0.12,-0.3)}, anchor=north, legend columns=-1, /tikz/every even column/.append style={column sep=8pt}}]
\nextgroupplot[ylabel={val. loss ($\downarrow$ better)}, ylabel style={align=center}, title={ViT-B}, xticklabels={}]
\addplot[cFT, solid, thick, mark=square*, mark size=1.8pt, forget plot, error bars/.cd, y dir=both, y explicit] coordinates {(3e-06,0.3681) +- (0,0.0122) (1e-05,0.3430) +- (0,0.0122) (3e-05,0.3198) +- (0,0.0135) (0.0001,0.3488) +- (0,0.0156)};
\addplot[cAccent, only marks, mark=o, mark size=4.5pt, very thick, forget plot] coordinates {(3e-05,0.3198)};
\addplot[cAttn, densely dashed, thick, mark=*, mark size=1.8pt, forget plot, error bars/.cd, y dir=both, y explicit] coordinates {(3e-05,0.3032) +- (0,0.0144) (0.0001,0.2931) +- (0,0.0157) (0.0003,0.2906) +- (0,0.0115) (0.001,0.3266) +- (0,0.0134)};
\addplot[cAccent, only marks, mark=o, mark size=4.5pt, very thick, forget plot] coordinates {(0.0001,0.2931)};
\nextgroupplot[title={Swin-B}, xticklabels={}]
\addplot[cFT, solid, thick, mark=square*, mark size=1.8pt, forget plot, error bars/.cd, y dir=both, y explicit] coordinates {(3e-06,0.3040) +- (0,0.0188) (1e-05,0.2912) +- (0,0.0155) (3e-05,0.3069) +- (0,0.0203) (0.0001,0.3827) +- (0,0.0066)};
\addplot[cAccent, only marks, mark=o, mark size=4.5pt, very thick, forget plot] coordinates {(1e-05,0.2912)};
\addplot[cAttn, densely dashed, thick, mark=*, mark size=1.8pt, forget plot, error bars/.cd, y dir=both, y explicit] coordinates {(3e-05,0.2714) +- (0,0.0154) (0.0001,0.2714) +- (0,0.0115) (0.0003,0.2833) +- (0,0.0114) (0.001,0.3111) +- (0,0.0109)};
\addplot[cAccent, only marks, mark=o, mark size=4.5pt, very thick, forget plot] coordinates {(0.0001,0.2714)};
\nextgroupplot[ylabel={test acc., \% ($\uparrow$ better)}, ylabel style={align=center}, xlabel={learning rate}]
\addplot[cFT, solid, thick, mark=square*, mark size=1.8pt, forget plot, error bars/.cd, y dir=both, y explicit] coordinates {(3e-06,90.1067) +- (0,0.1790) (1e-05,90.8800) +- (0,0.1929) (3e-05,91.3133) +- (0,0.1721) (0.0001,90.1700) +- (0,0.1646)};
\addplot[cAccent, only marks, mark=o, mark size=4.5pt, very thick, forget plot] coordinates {(3e-05,91.3133)};
\addplot[cAttn, densely dashed, thick, mark=*, mark size=1.8pt, forget plot, error bars/.cd, y dir=both, y explicit] coordinates {(3e-05,91.0067) +- (0,0.1845) (0.0001,91.0867) +- (0,0.1041) (0.0003,91.3933) +- (0,0.1474) (0.001,90.3600) +- (0,0.0866)};
\addplot[cAccent, only marks, mark=o, mark size=4.5pt, very thick, forget plot] coordinates {(0.0001,91.0867)};
\nextgroupplot[xlabel={learning rate}]
\addplot[cFT, solid, thick, mark=square*, mark size=1.8pt, forget plot, error bars/.cd, y dir=both, y explicit] coordinates {(3e-06,90.8900) +- (0,0.2751) (1e-05,91.3267) +- (0,0.3035) (3e-05,91.1567) +- (0,0.2532) (0.0001,88.5333) +- (0,0.0351)};
\addplot[cAccent, only marks, mark=o, mark size=4.5pt, very thick, forget plot] coordinates {(1e-05,91.3267)};
\addplot[cAttn, densely dashed, thick, mark=*, mark size=1.8pt, forget plot, error bars/.cd, y dir=both, y explicit] coordinates {(3e-05,91.6733) +- (0,0.2902) (0.0001,91.8300) +- (0,0.1473) (0.0003,91.5533) +- (0,0.2312) (0.001,90.6567) +- (0,0.1060)};
\addplot[cAccent, only marks, mark=o, mark size=4.5pt, very thick, forget plot] coordinates {(0.0001,91.8300)};
\addlegendimage{cFT, solid, thick, mark=square*, mark size=1.8pt}
\addlegendentry{full fine-tuning}
\addlegendimage{cAttn, densely dashed, thick, mark=*, mark size=1.8pt}
\addlegendentry{LoRA (attention)}
\addlegendimage{cAccent, only marks, mark=o, mark size=3.5pt, very thick}
\addlegendentry{rate used}
\addlegendimage{cGrey, mark=|, mark size=3pt}
\addlegendentry{$\pm$ std, 3 seeds}
\end{groupplot}
\end{tikzpicture}

\caption{Validation loss (top) and test accuracy (bottom) against the learning rate, 3 seeds. Circles mark the
rates used in the main comparison: for full fine-tuning the rate with the lowest validation loss, for LoRA $10^{-4}$.}
\label{fig:lr}
\end{figure}

\paragraph{Full fine-tuning.} On both models the validation loss has a clear minimum inside the range tried: at
$3\times10^{-5}$ on ViT-B and at $10^{-5}$ on Swin-B, with higher loss at the neighbouring rates on both sides. The
choice matters. On ViT-B, full fine-tuning at $10^{-5}$ is 0.43 $\pm$ 0.06~pp less accurate than at $3\times10^{-5}$
and at $10^{-4}$ it is 1.14 $\pm$ 0.19~pp less accurate (paired, all three seeds); with all five seeds the difference
between $3\times10^{-5}$ and $10^{-5}$ is 0.60 $\pm$ 0.23~pp, in every seed. On Swin-B, full fine-tuning at $10^{-4}$
is 2.8~pp below its best rate.

\paragraph{LoRA.} On ViT-B, $10^{-4}$ and $3\times10^{-4}$ have almost the same validation loss (difference 0.002
$\pm$ 0.005, both signs), and the two outer rates are worse in every seed. On Swin-B, $3\times10^{-5}$ and
$10^{-4}$ have the same validation loss (0.000 $\pm$ 0.005), and the two higher rates are worse in every seed. The
rate $10^{-4}$ is therefore at, or tied with, LoRA's best rate on both models, and it is the rate used for all LoRA
runs below.

\paragraph{Reference.} In the rest of the report, full fine-tuning means the rate $3\times10^{-5}$ on ViT-B and
$10^{-5}$ on Swin-B, each run with all five seeds.

\section{Accuracy of LoRA against full fine-tuning}
\label{sec:acc}

\begin{table}[!htb]
\centering\footnotesize
\setlength{\tabcolsep}{4pt}
\begin{tabular}{lcccccc}
\toprule
 & test acc. & vs full FT & seeds & val. & vs full FT, & best / \\
config & (\%) & (pp, paired) & above FT & loss & val.\ loss (paired) & stop ep. \\
\midrule
\rawinput{data/tab_acc.tex}
\bottomrule
\end{tabular}
\caption{LoRA (rank 8, all blocks, learning rate $10^{-4}$) and full fine-tuning at its best rate; mean $\pm$
standard deviation over 5 seeds. vs full FT = paired difference to full fine-tuning with the same seed, for test
accuracy and for validation loss (negative = LoRA has the lower loss). Seeds above FT = number of seeds in which
LoRA has the higher test accuracy. Best / stop ep.\ = mean best and stopping epoch. Per seed:
Appendix~\ref{app:seeds}.}
\label{tab:acc}
\end{table}

\begin{figure}[!htb]
\centering
% generated by scripts/make_figures.py; do not edit
\begin{tikzpicture}
\begin{groupplot}[group style={group size=2 by 1, horizontal sep=\figSep, vertical sep=1.3cm}, width=\figW, height=\figH, repaxis, xmin=0.4, xmax=4.6, xtick={1,2,3,4}, xticklabels={Full FT,Attn,MLP,Attn+MLP}, ymajorgrids, legend style={font=\scriptsize, draw=none, fill=none, at={(-0.12,-0.16)}, anchor=north, legend columns=-1, /tikz/every even column/.append style={column sep=8pt}}]
\nextgroupplot[title={ViT-B}, ylabel={test accuracy (\%)}]
\addplot[only marks, mark=square*, mark size=1.8pt, cFT, forget plot] coordinates {(0.82,91.51) (0.91,91.19) (1.00,91.24) (1.09,91.42) (1.18,91.19)};
\draw[cFT, line width=1.2pt] (axis cs:0.7,91.310) -- (axis cs:1.3,91.310);
\addplot[only marks, mark=*, mark size=1.8pt, cAttn, forget plot] coordinates {(1.82,91.17) (1.91,91.12) (2.00,90.97) (2.09,91.12) (2.18,91.13)};
\draw[cAttn, line width=1.2pt] (axis cs:1.7,91.102) -- (axis cs:2.3,91.102);
\addplot[only marks, mark=triangle*, mark size=1.8pt, cMLP, forget plot] coordinates {(2.82,91.30) (2.91,91.16) (3.00,91.33) (3.09,91.40) (3.18,91.57)};
\draw[cMLP, line width=1.2pt] (axis cs:2.7,91.352) -- (axis cs:3.3,91.352);
\addplot[only marks, mark=diamond*, mark size=1.8pt, cBoth, forget plot] coordinates {(3.82,91.44) (3.91,91.38) (4.00,91.57) (4.09,91.30) (4.18,91.47)};
\draw[cBoth, line width=1.2pt] (axis cs:3.7,91.432) -- (axis cs:4.3,91.432);
\draw[cFT, densely dashed] (axis cs:0.4,91.310) -- (axis cs:4.6,91.310);
\nextgroupplot[title={Swin-B}]
\addplot[only marks, mark=square*, mark size=1.8pt, cFT, forget plot] coordinates {(0.82,91.00) (0.91,91.60) (1.00,91.38) (1.09,91.19) (1.18,91.15)};
\draw[cFT, line width=1.2pt] (axis cs:0.7,91.264) -- (axis cs:1.3,91.264);
\addplot[only marks, mark=*, mark size=1.8pt, cAttn, forget plot] coordinates {(1.82,91.67) (1.91,91.96) (2.00,91.86) (2.09,92.00) (2.18,91.23)};
\draw[cAttn, line width=1.2pt] (axis cs:1.7,91.744) -- (axis cs:2.3,91.744);
\addplot[only marks, mark=triangle*, mark size=1.8pt, cMLP, forget plot] coordinates {(2.82,91.74) (2.91,91.64) (3.00,91.75) (3.09,91.94) (3.18,91.40)};
\draw[cMLP, line width=1.2pt] (axis cs:2.7,91.694) -- (axis cs:3.3,91.694);
\addplot[only marks, mark=diamond*, mark size=1.8pt, cBoth, forget plot] coordinates {(3.82,91.84) (3.91,91.54) (4.00,91.67) (4.09,91.99) (4.18,91.46)};
\draw[cBoth, line width=1.2pt] (axis cs:3.7,91.700) -- (axis cs:4.3,91.700);
\draw[cFT, densely dashed] (axis cs:0.4,91.264) -- (axis cs:4.6,91.264);
\addlegendimage{only marks, mark=*, mark size=1.8pt, cGrey}
\addlegendentry{one seed}
\addlegendimage{cGrey, line width=1.2pt}
\addlegendentry{mean of 5 seeds}
\addlegendimage{cFT, densely dashed}
\addlegendentry{mean of full FT}
\end{groupplot}
\end{tikzpicture}

\caption{Test accuracy of full fine-tuning and of the three LoRA targets, per seed.}
\label{fig:acc_seeds}
\end{figure}

\paragraph{Swin-B: LoRA is more accurate.} All three LoRA configurations are above full fine-tuning, by
0.43--0.48~pp, in 5, 5 and 4 of the 5 seeds (Table~\ref{tab:acc}, Figure~\ref{fig:acc_seeds}). All three
differences are detectable. In the learning-rate check, attention LoRA is above full fine-tuning at its best rate in
all three seeds, at both of its tied rates: by 0.50 $\pm$ 0.16~pp at $10^{-4}$ and by 0.35 $\pm$ 0.27~pp at
$3\times10^{-5}$.

\paragraph{ViT-B: LoRA is level with or slightly below full fine-tuning, and never ahead.} MLP and Attn+MLP do not differ detectably from
full fine-tuning ($+$0.04 $\pm$ 0.22 and $+$0.12 $\pm$ 0.21~pp). Attention LoRA is 0.21 $\pm$ 0.13~pp \emph{below}
it, in all five seeds, which is a detectable difference. Two further measurements show how far this difference can
be trusted:
\begin{itemize}
  \item At the rate $3\times10^{-4}$, which has the same validation loss as $10^{-4}$, attention LoRA is level with
        full fine-tuning: $+$0.08 $\pm$ 0.27~pp (3 seeds, Table~\ref{tab:lr}).
  \item The full fine-tuning reference was trained in a later set of runs than the LoRA runs of
        Table~\ref{tab:acc}. Attention LoRA at $10^{-4}$ was therefore trained again together with it, with three
        seeds. This repeated set is 0.19~pp more accurate than the runs of Table~\ref{tab:acc} on the same seeds,
        and it is level with full fine-tuning ($-$0.04 $\pm$ 0.29~pp). Two sets of ViT-B runs with identical
        settings and seeds can thus differ by about as much as the difference in question
        (Appendix~\ref{app:impl}).
\end{itemize}
We conclude that on ViT-B, LoRA is within about 0.2~pp of full fine-tuning: MLP and Attn+MLP show no detectable
difference, attention LoRA is at most 0.2~pp below, and there is no evidence that any of them is better.

\paragraph{Validation loss.} LoRA has the lower validation loss than full fine-tuning in all 30 seed pairs, by about
0.03 on ViT-B and 0.02 on Swin-B. On Swin-B this agrees with the test accuracy. On ViT-B it does not turn into
higher test accuracy; a lower loss does not have to mean more correct predictions.

\paragraph{What this comparison depends on.} Against full fine-tuning at $10^{-5}$ instead of $3\times10^{-5}$, LoRA
would appear 0.4--0.7~pp better on ViT-B in every seed. The outcome of a comparison between LoRA and full
fine-tuning can therefore be decided by how well full fine-tuning is tuned. Three further limits apply: the rates
were varied with three seeds and for the attention target only; accuracy on CIFAR-100 is close to saturated for
these models; and with five seeds, differences of about 0.2~pp can be neither shown nor ruled out. We did not test
why LoRA, which trains less than 2\% of the weights, loses nothing on Swin-B.

\section{Targets and rank}
\label{sec:targets}

\paragraph{Targets.} On Swin-B no difference between the three targets is detectable (paired differences between
$-$0.05 and $+$0.01~pp). On ViT-B at learning rate $10^{-4}$, MLP and Attn+MLP are more accurate than attention
alone, by 0.25 $\pm$ 0.16 and 0.33 $\pm$ 0.16~pp, in every seed. This ranking holds at that rate only: the rate was
varied for attention alone, and at $3\times10^{-4}$ attention LoRA is level with the other two at $10^{-4}$
($+$0.13 $\pm$ 0.24~pp against MLP, $-$0.07 $\pm$ 0.24~pp against Attn+MLP; 3 seeds). We therefore do not conclude
that one target is more accurate than another. What differs clearly between the targets is their cost
(Section~\ref{sec:mem}).

\paragraph{Rank.}
\label{sec:rank}
ViT-B was also trained with rank 16 for all three targets, with five seeds, in an earlier set of runs that includes
its own rank-8 runs. Rank 16 gives no gain: the paired differences to rank 8 are $-$0.15 $\pm$ 0.12~pp (attention),
$-$0.02 $\pm$ 0.13~pp (MLP) and $+$0.05 $\pm$ 0.12~pp (Attn+MLP). For attention, rank 16 is detectably lower in the
library implementation (in all five seeds), but our own implementation shows no such effect ($+$0.02 $\pm$ 0.04~pp). Peak memory rises by less than 0.04~GiB
(Table~\ref{tab:rank} in the appendix). Rank 16 was not run on Swin-B.

\section{Training memory}
\label{sec:mem}

\begin{figure}[!htb]
\centering
% generated by scripts/make_figures.py; do not edit
\begin{tikzpicture}
\begin{groupplot}[group style={group size=2 by 1, horizontal sep=\figSep, vertical sep=1.3cm}, width=\figW, height=\figH, repaxis, ybar, /pgf/bar width=14pt, xmin=0.4, xmax=4.6, ymin=0, xtick={1,2,3,4}, xticklabels={Full FT,Attn,MLP,Attn+MLP}, ymajorgrids, enlarge y limits={upper, value=0.22}]
\nextgroupplot[title={ViT-B}, ylabel={peak training memory (GiB)}]
\addplot[fill=cFT, draw=black, bar shift=0pt] coordinates {(1,8.174)};
\node[font=\scriptsize, anchor=south, align=center] at (axis cs:1,8.174) {8.2 GiB};
\addplot[fill=cAttn, draw=black, bar shift=0pt] coordinates {(2,5.459)};
\node[font=\scriptsize, anchor=south, align=center] at (axis cs:2,5.459) {5.5 GiB\\$-$33\%};
\addplot[fill=cMLP, draw=black, bar shift=0pt] coordinates {(3,7.011)};
\node[font=\scriptsize, anchor=south, align=center] at (axis cs:3,7.011) {7.0 GiB\\$-$14\%};
\addplot[fill=cBoth, draw=black, bar shift=0pt] coordinates {(4,7.643)};
\node[font=\scriptsize, anchor=south, align=center] at (axis cs:4,7.643) {7.6 GiB\\$-$6\%};
\nextgroupplot[title={Swin-B}]
\addplot[fill=cFT, draw=black, bar shift=0pt] coordinates {(1,14.121)};
\node[font=\scriptsize, anchor=south, align=center] at (axis cs:1,14.121) {14.1 GiB};
\addplot[fill=cAttn, draw=black, bar shift=0pt] coordinates {(2,9.549)};
\node[font=\scriptsize, anchor=south, align=center] at (axis cs:2,9.549) {9.5 GiB\\$-$32\%};
\addplot[fill=cMLP, draw=black, bar shift=0pt] coordinates {(3,12.007)};
\node[font=\scriptsize, anchor=south, align=center] at (axis cs:3,12.007) {12.0 GiB\\$-$15\%};
\addplot[fill=cBoth, draw=black, bar shift=0pt] coordinates {(4,13.290)};
\node[font=\scriptsize, anchor=south, align=center] at (axis cs:4,13.290) {13.3 GiB\\$-$6\%};
\end{groupplot}
\end{tikzpicture}

\caption{Peak training memory at batch size 64, with the saving against full fine-tuning. The values are
identical in all 5 seeds.}
\label{fig:mem_targets}
\end{figure}

The pattern is the same on both models (Figure~\ref{fig:mem_targets}; Table~\ref{tab:mem} in the appendix). LoRA
reduces the static memory from about 1.3 to 0.33--0.35~GiB, because gradients and optimiser state are kept only for
the adapters. But static memory is the small part: the peak of full fine-tuning is 8.2~GiB on ViT-B and 14.1~GiB on Swin-B, of
which only 1.3~GiB is static; the rest is activation memory. How much of it LoRA removes depends on the target, for the reason given in
Section~\ref{sec:bg_memory}: a frozen layer without an adapter does not keep its input, an adapted layer does, and
the input of the second MLP layer is four times wider than that of the attention projections. On both models
attention LoRA saves about a third of the peak memory (33\% and 32\%), MLP LoRA about a seventh (14\% and 15\%), and
Attn+MLP about 6\%.

\begin{figure}[!htb]
\centering
% generated by scripts/make_figures.py; do not edit
\begin{tikzpicture}
\begin{groupplot}[group style={group size=2 by 1, horizontal sep=\figSep, vertical sep=1.3cm}, width=\figW, height=\figH, repaxis, xmode=log, log basis x=2, xtick={1,2,4,8,16,32,64,128}, xticklabels={1,2,4,8,16,32,64,128}, ymin=0, ymax=80, ymajorgrids, legend style={font=\footnotesize, draw=none, fill=none}, legend pos=north east]
\nextgroupplot[title={ViT-B}, xlabel={batch size}, ylabel={saving against full FT (\%)}]
\addplot[cAttn, solid, thick, mark=*, mark size=1.8pt] coordinates {(1,73.8) (2,68.0) (4,59.3) (8,48.1) (16,42.0) (32,36.9) (64,33.2) (128,31.0)};
\addplot[cMLP, densely dashed, thick, mark=triangle*, mark size=1.8pt] coordinates {(1,72.1) (2,64.8) (4,53.1) (8,37.7) (16,28.0) (32,19.9) (64,14.2) (128,10.9)};
\addplot[cBoth, densely dotted, thick, mark=diamond*, mark size=1.8pt] coordinates {(1,71.2) (2,63.2) (4,50.4) (8,33.3) (16,22.2) (32,13.0) (64,6.5) (128,2.6)};
\nextgroupplot[title={Swin-B}, xlabel={batch size}]
\addplot[cAttn, solid, thick, mark=*, mark size=1.8pt] coordinates {(1,69.8) (2,60.8) (4,51.2) (8,44.2) (16,38.2) (32,34.4) (64,32.4) (128,31.2)};
\addlegendentry{Attn}
\addplot[cMLP, densely dashed, thick, mark=triangle*, mark size=1.8pt] coordinates {(1,67.1) (2,55.9) (4,42.8) (8,32.8) (16,23.8) (32,18.1) (64,15.0) (128,13.2)};
\addlegendentry{MLP}
\addplot[cBoth, densely dotted, thick, mark=diamond*, mark size=1.8pt] coordinates {(1,65.5) (2,53.0) (4,38.4) (8,26.7) (16,16.1) (32,9.5) (64,5.9) (128,3.8)};
\addlegendentry{Attn+MLP}
\end{groupplot}
\end{tikzpicture}

\caption{Saving in peak training memory against full fine-tuning for batch sizes 1 to 128 (one training step,
profiler). Full fine-tuning of Swin-B does not fit on the 44.25~GiB partition at batch size 256.}
\label{fig:saving_batch}
\end{figure}

The relative saving is largest for small batches (Figure~\ref{fig:saving_batch}; Table~\ref{tab:batch} in the
appendix). At batch size 1 every LoRA target needs 66--74\% less memory than full fine-tuning, because the static
memory, which LoRA cuts to a quarter, is then the larger part. As the batch grows, activation memory takes over and
the saving falls towards the share of the activations that the target avoids: at batch size 128 it is 31\% for
attention on both models, 11--13\% for MLP and 3--4\% for Attn+MLP. Small batches are where memory-limited devices
operate, so this is the favourable case for LoRA.

\section{Training time}
\label{sec:time}

\begin{figure}[!htb]
\centering
% generated by scripts/make_figures.py; do not edit
\begin{tikzpicture}
\begin{groupplot}[group style={group size=2 by 1, horizontal sep=\figSep, vertical sep=1.3cm}, width=\figW, height=\figH, repaxis, ybar stacked, /pgf/bar width=14pt, xmin=0.4, xmax=4.6, ymin=0, xtick={1,2,3,4}, xticklabels={Full FT,Attn,MLP,Attn+MLP}, ymajorgrids, legend style={font=\footnotesize, draw=none, fill=none}, legend pos=north east, legend columns=-1, enlarge y limits={upper, value=0.22}]
\nextgroupplot[title={ViT-B}, ylabel={time per training batch (ms)}]
\addplot[fill=cGrey!35, draw=black] coordinates {(1,208.6) (2,219.8) (3,224.9) (4,235.5)};
\addplot[fill=cAccent, draw=black] coordinates {(1,407.9) (2,236.9) (3,235.1) (4,264.0)};
\addplot[fill=black, draw=black] coordinates {(1,5.5) (2,1.1) (3,1.0) (4,1.6)};
\nextgroupplot[title={Swin-B}]
\addplot[fill=cGrey!35, draw=black] coordinates {(1,245.8) (2,263.4) (3,273.4) (4,289.9)};
\addlegendentry{forward}
\addplot[fill=cAccent, draw=black] coordinates {(1,435.9) (2,271.1) (3,271.6) (4,315.4)};
\addlegendentry{backward}
\addplot[fill=black, draw=black] coordinates {(1,6.5) (2,1.4) (3,1.1) (4,2.2)};
\addlegendentry{optimiser}
\end{groupplot}
\end{tikzpicture}

\caption{Time per training batch of 64 images, split into forward pass, backward pass and optimiser step (mean
over 5 seeds).}
\label{fig:time_batch}
\end{figure}

A training epoch is 19--26\% shorter with LoRA on ViT-B (323--354~s against 439~s) and 12--22\% shorter on Swin-B
(378--428~s against 485~s); the standard deviation over seeds is about 1~s or less (Table~\ref{tab:time} in the
appendix).
Figure~\ref{fig:time_batch} shows where the time goes. The backward pass becomes much faster (from 408 to
235--264~ms per batch on ViT-B, from 436 to 271--315~ms on Swin-B), because no gradients are computed for the frozen
weights, and the optimiser step, which updates 53--165 times fewer values, becomes negligible. The forward pass
becomes slightly \emph{slower}, because during training the adapters are separate layers that add operations. For
the same reason a validation epoch, which consists of forward passes only, is 6--13\% slower with LoRA on ViT-B and
7--18\% slower on Swin-B; it is a small part of the total (17--23~s).

\section{Stored size}
\label{sec:size}

Counted as 4 bytes per trained value, a LoRA task with adapters on all blocks is stored in 2.0--4.8~MiB on ViT-B and 2.6--6.3~MiB on Swin-B,
against 328 and 331~MiB for a fully fine-tuned model. This is an advantage only when the base model is already on
the device, for example when several tasks share it and each adds its own adapter. For a single task nothing is
gained: the merged LoRA model has exactly the size of a fully fine-tuned one (Chapter~\ref{ch:deploy}).

\section{Summary}

With both methods at their best learning rates, LoRA at rank 8 is more accurate than full fine-tuning on Swin-B and
level with or slightly below it on ViT-B. The cost differs clearly between the targets, in the same way on both models: attention LoRA
saves a third of the training memory at batch size 64, and more at smaller batches, and shortens the training epoch
by about a quarter; adapters on the MLP save much less memory. No target is reliably more accurate than another, so
attention is the target with the best balance of accuracy and cost, and the following chapters use it.
